This chapter demonstrates the application of Omnisoot in various use-cases. Each section of the chapter represents a paper, in which Omnisoot was used for predicting soot yield, morphology and size distribution. The uncertainties from gas-phase chemistry and empirically-derived inception pathways were highlighted and discussed. The carbon flux from gas to solid via inception, HACA and PAH adsorption were investigated to understand their effect on soot properties. A special attention was paid to methane pyrolysis in shock-tubes with different compositions, and temperatures, and pressures because temperature and pressure ranges achieved in shock-tube experiments can reach 2300~K and 5~atm that are relevant to CB synthesis in industrial plasma reactors~\citep{fulcheri2023energy}, which is a major motivations for developing Omnisoot. Additionally, a unique experimental data-set were compiled and processed in a joint program between Carleton-Monolith-Stanford. The shock-tube experimental campaign cover a wide range of high methane loadings (5\%, 10\%, and 30\%) and temperature sweep of 1800~K-2500~K, and it provides simultaneous species and particulate measurements. 


%A special attention is paid to the simulation of methane pyrolysis in shock-tubes with different compositions, and temperatures, and pressures because temperature and pressure ranges achieved in shock-tube experiments can reach 2300~K and 10~atm that are relevant to industrial plasma reactors~\citep{fulcheri2023energy}, so it can inform the process design of CB synthesis, which is one the main motivations for development of Omnisoot. Moreover, the data acquisition was performed for both species and soot by Hanson Group in Stanford in collaboration with Carleton and Monolith, the industrial collaborator of the research. The time-resolved species and soot measurements makes the shock-tube data unique for analyzing the carbon flux from methane to soot, and explaining the uncertainties in prediction of soot yield and morphology that originate from gas-phase chemistry, and inception and surface growth rates. The performance of the implemented inception models are evaluated in terms of their temperature-dependence and the balance of inception and PAH adsorption that determines primary particle diameter.

%Shock tubes provide valuable information on the initial stages of soot formation, such as inception and surface growth, due to short residence times of the process ($\approx1$-5~ms). On the other hand, flow reactors have longer residence times (0.1-1~s) where coagulation also becomes important in determining soot morphology. These conditions are more relevant to industrial reactors~\citep{fulcheri2023energy}, but the species~\citep{Norinaga2008} and soot measurements~\citep{araki2021effects} are only available at the end of reactor, so it is difficult to get a full picture of evolution of soot particles along the reactor. We use Omnisoot to study soot formation during ethylene pyrolysis in two different laminar reactor cases. We will also simulate a PSR connected to a PFR to evaluate the capability of Omnisoot in predicting soot yield, morphology and PSD.